\subsection{本质素数谱、记录轴最小性与哥德尔化外置维数律}\label{subsec:conclusion-essential-prime-spectrum-primorial-witt}
本小节在既有素数寄存器框架
\[
\mathbf P=\NN^{(\mathbb P)},
\qquad
N(\mathbf r)=\prod_{p\in\mathbb P}p^{\mathbf r(p)}
\]
下，给出分歧素数谱的记录轴最小性判据，并将其与 primorial 维数律及 Big--Witt 指纹的哥德尔化可乘结构统一为同一外置语义。

\begin{definition}[分歧素数谱与本质素数谱]\label{def:conclusion-essential-prime-spectrum}
设 \(\Phi(\lambda,y)\in\QQ[\lambda,y]\) 绝对不可约，令
\[
D_0:=\mathrm{den}(\Phi),\qquad
\widetilde\Phi:=D_0\Phi\in\ZZ[\lambda,y]
\]
为本原整系数模型。设 \(C\) 为 \(\widetilde\Phi(\lambda,y)=0\) 的正规化光滑射影模型，并记有限映射
\[
\pi_y:C\longrightarrow \PP^1_y.
\]
定义
\[
\Delta(y):=\mathrm{Disc}_{\lambda}\bigl(\widetilde\Phi(\lambda,y)\bigr)\in\ZZ[y],
\qquad
\mathfrak D:=\mathrm{Disc}_{y}\bigl(\Delta(y)\bigr)\in\ZZ\setminus\{0\},
\]
以及
\[
S_{\mathrm{ess}}(\pi_y):=\{\,p\in\mathbb P:\ p\mid D_0\mathfrak D\,\}.
\]
称 \(S_{\mathrm{ess}}(\pi_y)\) 为该覆盖在稳定反演语义下的本质素数谱。
\end{definition}

\begin{theorem}[分歧素数谱的记录轴最小性]\label{thm:conclusion-essential-prime-axis-minimality}
\leanverified{paper\_conclusion\_essential\_prime\_axis\_minimality}
沿用定义 \ref{def:conclusion-essential-prime-spectrum}。
\begin{enumerate}
  \item \textbf{充分性（好素数的有限索引外置）}：
  若 \(p\notin S_{\mathrm{ess}}(\pi_y)\)，则在
  \[
  U_p:=\{\,y\in\PP^1(\overline{\QQ}_p):\ v_p(\Delta(y))=0\,\}
  \]
  上，存在有限集合 \(A_p\)（\(|A_p|\le \deg(\pi_y)\)）及映射
  \[
  \rho_p:\pi_y^{-1}(U_p)\longrightarrow A_p
  \]
  使得
  \[
  P\longmapsto \bigl(\pi_y(P),\rho_p(P)\bigr)
  \]
  在 \(\pi_y^{-1}(U_p)\) 上注入，且该标号与所有 \(p^k\)-级 Hensel 提升相容。

  \item \textbf{必要性（坏素数的无限记录轴不可压缩）}：
  若 \(p\in S_{\mathrm{ess}}(\pi_y)\)，则不存在有限集合 \(A\) 与映射
  \[
  \rho:\pi_y^{-1}(V)\to A
  \]
  能在“趋近坏特化点的全部 \(p\)-进同余塔”上同时保持注入与提升稳定，其中 \(V\subset\PP^1(\overline{\QQ}_p)\) 为任意与分支集相交可任意逼近的穿孔邻域族。
  等价地，若要求在该坏素数方向上实现稳定可逆外置，则必须保留无限精度的 \(p\)-进记录分量（在完备化中等价于保留 \(\ZZ_p\) 轴）。
\end{enumerate}
\end{theorem}
\begin{proof}
对 \(p\notin S_{\mathrm{ess}}(\pi_y)\)，有 \(p\nmid D_0\) 与 \(p\nmid\mathfrak D\)。
前者给出良好整模型，后者保证 \(\Delta(y)\) 在模 \(p\) 下无重根，故在 \(U_p\) 上纤维对应的 \(\lambda\)-根均为单根且分支不碰撞。
于是每个模 \(p\) 单根由 Hensel 引理唯一提升为 \(p^k\)-根塔；以根的模 \(p\) 类别作索引即可得到有限集合 \(A_p\) 与注入标号 \(\rho_p\)，并自动满足提升稳定。

对 \(p\in S_{\mathrm{ess}}(\pi_y)\)，若 \(p\mid D_0\) 则整模型在该素数方向发生退化；若 \(p\mid\mathfrak D\) 则分支点在特征 \(p\) 下碰撞。
两种情形均产生可任意高阶贴近的不同分支提升：存在趋近坏特化的 \(y_k\) 及不同点 \(P_k\neq Q_k\) 使其在 \(p^k\)-层面不可由任何固定有限记号系统稳定分离。
因此有限 \(A\) 无法在所有同余层上同时维持注入与提升相容；若要求稳定反演，必须引入无限 \(p\)-进记录轴。
证毕。
\end{proof}

\begin{corollary}[Lee--Yang 四叶覆盖的奇素本质谱]\label{cor:conclusion-lee-yang-essential-odd-primes}
设 Lee--Yang 四叶覆盖由 \(\Pi(\lambda,y)=0\) 给出，并满足判别式关系
\[
\mathrm{Disc}_{\lambda}(4\Pi)=-y(y-1)\,\mathrm{PLY}(y),
\qquad
\mathrm{Disc}\bigl(\mathrm{PLY}\bigr)=-3^9\cdot 31^2\cdot 37.
\]
又其归一化双有理模型可写为
\[
E:\ w^2=4\lambda^3-4\lambda+1,
\qquad
\Delta_E=37.
\]
则在“\(p\)-进提升稳定反演外置”语义下，其奇素本质谱满足
\[
\{3,31,37\}\subseteq S_{\mathrm{ess}}(\pi_y).
\]
因此 \(37\)（坏约化素数）与 \(3,31\)（分支碰撞素数）均不可由有限索引稳定消去。
\end{corollary}
\begin{proof}
由给定判别式分解，\(\mathrm{PLY}\) 的判别式素支撑含 \(\{3,31,37\}\)，相应分支退化必进入
\(
\mathrm{Disc}_y(\Delta)
\)
的坏素集合。
同时 \(\Delta_E=37\) 给出结构性坏素数 \(37\)。
代入定义 \ref{def:conclusion-essential-prime-spectrum} 与定理 \ref{thm:conclusion-essential-prime-axis-minimality} 即得结论。证毕。
\end{proof}

\begin{theorem}[Lee--Yang 四次覆盖的本质素数谱]\label{thm:conclusion-leyang-quartic-essential-prime-spectrum-explicit}
令 \(\Pi(\lambda,y)\in\ZZ[\lambda,y]\) 为定理 \ref{thm:xi-terminal-zm-leyang-elliptic-structure} 的谱四次，
并令 \(\pi_y:C\to\PP^1_y\) 为平面曲线 \(\Pi(\lambda,y)=0\) 的规范化所诱导的次数 \(4\) 覆盖。
则在定义 \ref{def:conclusion-essential-prime-spectrum} 的口径下，其本质素数谱满足
\[
\boxed{\ S_{\mathrm{ess}}(\pi_y)=\{2,3,31,37\}\ }.
\]
\end{theorem}
\begin{proof}
由于 \(\Pi\) 系数整，故定义 \ref{def:conclusion-essential-prime-spectrum} 中 \(D_0=\mathrm{den}(\Pi)=1\)。
另一方面，由结论 \ref{con:xi-terminal-zm-discriminant-leyang} 的判别式恒等式
\[
\Disc_{\lambda}(4\Pi)=-y(y-1)P_{\mathrm{LY}}(y),
\]
并结合结论 \ref{con:xi-terminal-zm-leyang-discriminant-ridge-polydisc-badprimes} 的显式判别式分解
\[
\Disc_y\!\bigl(-y(y-1)P_{\mathrm{LY}}(y)\bigr)=-2^{20}\cdot 3^{15}\cdot 31^{2}\cdot 37,
\]
可知 \(\mathfrak D=\Disc_y(\Disc_\lambda(\widetilde\Pi))\) 的素支撑恰为 \(\{2,3,31,37\}\)。
代回定义 \ref{def:conclusion-essential-prime-spectrum} 即得结论。证毕。
\end{proof}

\begin{theorem}[本质素数谱的 \(37\)-耦合分解]\label{thm:conclusion-leyang-37-coupling-decomposition}
在定理 \ref{thm:conclusion-leyang-quartic-essential-prime-spectrum-explicit} 的记号下，令 \(E_{\min}/\QQ\) 为定理 \ref{thm:xi-terminal-zm-leyang-elliptic-structure} 的极小椭圆曲线，
并记其二挠分裂域
\[
K_{E[2]}:=\QQ\bigl(E_{\min}[2]\bigr),
\]
以及 Lee--Yang 三次 \(P_{\mathrm{LY}}(y)\) 的分裂域
\[
K_{\mathrm{LY}}:=\mathrm{Spl}\bigl(P_{\mathrm{LY}}(y)\bigr).
\]
则
\[
\mathrm{Ram}(K_{E[2]})=\{2,37\},\qquad
\mathrm{Ram}(K_{\mathrm{LY}})=\{3,31,37\},
\]
并且
\[
S_{\mathrm{ess}}(\pi_y)=\mathrm{Ram}(K_{E[2]})\cup \mathrm{Ram}(K_{\mathrm{LY}}),
\qquad
\mathrm{Ram}(K_{E[2]})\cap \mathrm{Ram}(K_{\mathrm{LY}})=\{37\}.
\]
\end{theorem}
\begin{proof}
由定理 \ref{thm:xi-terminal-zm-leyang-elliptic-structure}，\(K_{E[2]}\) 为 \(E_{\min}\) 的二除法多项式 \(4x^3-4x+1\) 的分裂域，且
\[
\Disc(4x^3-4x+1)=2^4\cdot 37,
\]
故 \(\mathrm{Ram}(K_{E[2]})=\{2,37\}\)。
另一方面，由结论 \ref{con:xi-terminal-zm-leyang-branchpolynomial-discriminant-modp-collisions}，
\[
\Disc\bigl(P_{\mathrm{LY}}(y)\bigr)=-3^9\cdot 31^2\cdot 37,
\]
从而 \(\mathrm{Ram}(K_{\mathrm{LY}})=\{3,31,37\}\)。
两式并集与交集结论遂立即给出。
最后，定理 \ref{thm:conclusion-leyang-quartic-essential-prime-spectrum-explicit} 已确定 \(S_{\mathrm{ess}}(\pi_y)=\{2,3,31,37\}\)，而右端并集亦恰为该集合，故等式成立。证毕。
\end{proof}

\begin{proposition}[本原性禁止全局二次--二次中间覆盖]\label{prop:conclusion-leyang-no-global-quadratic-quadratic-factorization}
沿用定理 \ref{thm:conclusion-leyang-quartic-essential-prime-spectrum-explicit} 的记号。
则 \(\pi_y\) 不存在非平凡中间二次覆盖 \(C\to C'\to\PP^1_y\)（两段次数均为 \(2\)）。
等价地，不存在一种对所有良素数同时有效的全局“二比特外置”，将同一纤维内的四张纸稳定分解为两层彼此独立的二元决策。
\end{proposition}
\begin{proof}
由推论 \ref{cor:xi-terminal-zm-leyang-cover-primitive}，覆盖 \(\pi_y\) 为本原覆盖。
而次数 \(4\) 覆盖分解为两段二次覆盖当且仅当其在四张纸上保留一个 \(2+2\) 的分块结构，亦即覆盖不本原。
因此该分解不可能存在。证毕。
\end{proof}

\endinput



\begin{theorem}[分支值结构强制的环路证书秩下界]\label{thm:conclusion-leyang-loop-certificate-rank-lower-bound}
\leanpartial{paper\_conclusion\_leyang\_loop\_certificate\_rank\_lower\_bound}{The Lean signature is too weak as stated: an injective map \((\pi,\kappa): \Omega \to B \times A\) does not force \(|B|-1 \le \mathrm{finrank}_{\mathbb Z}(A)\) without additional loop-space or surjectivity hypotheses.}
设 \(p:X\to\PP^1\) 为有限分歧覆叠，分支值集合记为 \(B\subset\PP^1\)，并令 \(U:=\PP^1\setminus B\)。
考虑以端点作为可观测读出 \(\pi:\Omega\to U\) 的路径/解析延拓态空间 \(\Omega\)，并设存在一个离散阿贝尔群 \(A\) 与映射 \(\kappa:\Omega\to A\) 使得联合读出 \((\pi,\kappa)\) 在对象层为单射。
则必有
\[
\rank_{\ZZ}A\ \ge\ \rank_{\ZZ}H_1(U;\ZZ)=|B|-1.
\]
在 Lee--Yang 设定下，若覆盖的分支值集合为
\[
B_y=\{\infty,0,1,y_1,y_2,y_3\},
\]
其中 \(y_i\) 为多项式 \(\mathrm{PLY}(y)\) 的三根（参见推论 \ref{cor:conclusion-lee-yang-essential-odd-primes} 的判别式分解），则任何保持扩展保存语义的环路证书群都满足 \(\rank_{\ZZ}A\ge 5\)。
\end{theorem}

\begin{proof}
对去点球面 \(U=\PP^1\setminus B\)，有 \(\pi_1(U)\) 为自由群 \(F_{|B|-1}\)，其阿贝尔化同构于
\[
H_1(U;\ZZ)\cong \ZZ^{|B|-1}.
\]
若 \(\rank_{\ZZ}A<|B|-1\)，则任意群同态 \(H_1(U;\ZZ)\to A\) 均有非平凡核，因而存在非零同调类之环路在 \(A\) 中同像。
这意味着存在两条以同一端点为读出的解析延拓路径，其相差一个在证书中不可见的环路而仍在 \((\pi,\kappa)\) 下不可区分，与单射性矛盾。
故 \(\rank_{\ZZ}A\ge |B|-1\)。
Lee--Yang 情形下 \(|B_y|=6\)，代入即得 \(\rank_{\ZZ}A\ge 5\)。证毕。
\end{proof}

\endinput


\begin{proposition}[primorial 最优策略下的最小素数轴数维数律]\label{prop:conclusion-primorial-axis-dimension-law}
设
\[
p_k^\#:=\prod_{i=1}^{k}p_i
\]
为第 \(k\) 个 primorial。给定容量门槛 \(P_m(\varepsilon)>1\)，定义
\[
k_{\min}(m,\varepsilon):=
\min\{\,k\ge 1:\ p_k^\#>P_m(\varepsilon)\,\}.
\]
若 \(P_m(\varepsilon)\to\infty\)，则
\[
k_{\min}(m,\varepsilon)
\sim
\frac{\log P_m(\varepsilon)}{\log\log P_m(\varepsilon)}.
\]
特别地，当
\[
P_m(\varepsilon)=\exp(m\gamma),\qquad \gamma>0,
\]
有
\[
k_{\min}(m,\varepsilon)
\sim
\frac{m\gamma}{\log(m\gamma)}.
\]
因此在最优互异素数模数策略下，外置所需素数轴数为 \(m/\log m\) 级次线性增长。
\end{proposition}
\begin{proof}
设 \(\vartheta(x):=\sum_{p\le x}\log p\)。则
\[
\log p_k^\#
=\sum_{i=1}^{k}\log p_i
=\vartheta(p_k).
\]
由素数定理，
\[
\vartheta(x)\sim x,\qquad p_k\sim k\log k,
\]
故
\[
\log p_k^\#\sim p_k\sim k\log k.
\]
令 \(L:=\log P_m(\varepsilon)\)，阈值条件 \(p_k^\#>P_m(\varepsilon)\) 等价于 \(k\log k\sim L\) 的反演问题，从而
\[
k\sim \frac{L}{\log L}
=\frac{\log P_m(\varepsilon)}{\log\log P_m(\varepsilon)}.
\]
取 \(L=m\gamma\) 即得第二式。证毕。
\end{proof}

\begin{theorem}[素数寄存器的精度--最坏模长--平方时间下界]\label{thm:conclusion-prime-register-worst-modulus-mixing-lower-bound}
\leanverified{paper\_conclusion\_prime\_register\_worst\_modulus\_mixing\_lower\_bound}
令 \(B_m(\varepsilon)>1\) 为固定窗口尺度 \(m\) 与碰撞预算 \(\varepsilon\) 所对应的容量阈值，并定义
\[
k^\ast(m;\varepsilon):=\min\left\{k\ge 1:\ p_k^\#:=\prod_{i=1}^k p_i\ge B_m(\varepsilon)\right\}.
\]
对任意有限素集 \(P\subset\mathbb P\)，考虑离散寄存器
\[
G_P:=\prod_{p\in P}\ZZ/p\ZZ,
\qquad
|G_P|=\prod_{p\in P}p,
\]
并假设其满足体积覆盖约束 \(|G_P|\ge B_m(\varepsilon)\)。
令
\[
L(P):=\max P,
\qquad
L_{\min}(m;\varepsilon):=\min\left\{L(P):\ \prod_{p\in P}p\ge B_m(\varepsilon)\right\}.
\]
则
\[
L_{\min}(m;\varepsilon)=p_{k^\ast(m;\varepsilon)}.
\]
进一步，令 \(W_P\) 为 \(G_P\) 上的惰性简单随机游走（每步以正概率对某一坐标加 \(\pm1\)，并以正概率保持不动），其谱隙 \(\Delta(W_P)\) 满足
\[
\Delta(W_P)\ \lesssim\ \frac{1}{L(P)^2},
\]
从而松弛时间 \(\tau_{\mathrm{rel}}(W_P):=\Delta(W_P)^{-1}\) 满足
\[
\tau_{\mathrm{rel}}(W_P)\ \gtrsim\ L(P)^2\ \ge\ L_{\min}(m;\varepsilon)^2\ =\ p_{k^\ast(m;\varepsilon)}^2.
\]
\end{theorem}

\begin{proof}
先证最小最坏模长结论。
对任意整数 \(L\ge 2\) 定义
\[
L^\#:=\prod_{p\le L}p.
\]
若 \(P\subseteq\{p:\ p\le L\}\)，则 \(\prod_{p\in P}p\le L^\#\)。
因此若存在 \(P\) 使 \(\prod_{p\in P}p\ge B_m(\varepsilon)\)，并令 \(L=L(P)\)，则必有 \(L^\#\ge B_m(\varepsilon)\)。
反之若 \(L^\#\ge B_m(\varepsilon)\)，取 \(P=\{p:\ p\le L\}\) 即满足约束。
故
\[
L_{\min}(m;\varepsilon)=\min\{L:\ L^\#\ge B_m(\varepsilon)\}.
\]
注意 \(L^\#\) 仅在 \(L\) 取素数处跳变，故该最小值必为某个素数 \(p_k\)，并且等价于
\(
\prod_{i=1}^k p_i=p_k^\#\ge B_m(\varepsilon)
\) 的最小 \(k\)，即 \(L_{\min}(m;\varepsilon)=p_{k^\ast(m;\varepsilon)}\)。

再证平方时间下界。
对循环群 \(\ZZ/p\ZZ\) 的惰性简单随机游走，其非平凡特征值由 \(\cos(2\pi j/p)\) 控制，从而谱隙满足
\(
\Delta\lesssim 1-\cos(2\pi/p)\asymp p^{-2}
\)。
对直积群 \(G_P\) 的坐标更新链，谱隙由最慢坐标（即最大模长 \(L(P)\)）控制，因此
\(
\Delta(W_P)\lesssim L(P)^{-2}
\)，从而 \(\tau_{\mathrm{rel}}(W_P)\gtrsim L(P)^2\)。
代入 \(L(P)\ge L_{\min}(m;\varepsilon)\) 得结论。证毕。
\end{proof}

\endinput


\begin{theorem}[二阶素数维度律：临界比特预算到 primorial 轴数的反演]\label{thm:conclusion-primorial-axis-second-order}
沿用推论 \ref{cor:pom-oracle-critical-window-bit-inversion} 的记号。对固定 \(p\in(0,1)\)，令 \(B_m(p)\) 为达到成功率 \(p+o(1)\) 的最小比特预算，满足
\[
B_m(p)
=
\beta_{\mathrm{crit}}m
+
\frac{\sigma}{\log 2}\Phi^{-1}(p)\sqrt m
+
o(\sqrt m).
\]
记
\[
P_k:=\prod_{i=1}^{k}p_i,
\qquad
r_m(p):=\min\{k\in\NN:\ \log_2 P_k\ge B_m(p)\},
\]
并定义自然对数尺度
\[
S_m(p):=(\log 2)\,B_m(p).
\]
则当 \(m\to\infty\) 时：
\begin{enumerate}
  \item
  \[
  r_m(p)=\frac{S_m(p)}{W(S_m(p))}\bigl(1+o(1)\bigr),
  \]
  其中 \(W\) 为 Lambert \(W\) 函数；
  \item 记 \(a:=(\log 2)\beta_{\mathrm{crit}},\ b:=\sigma\Phi^{-1}(p)\)，则
  \[
  r_m(p)
  =
  \frac{a}{\log m}\,m
  +
  \frac{b}{\log m}\,\sqrt m
  +
  O\!\left(\frac{m\log\log m}{(\log m)^2}\right)
  +
  o\!\left(\frac{\sqrt m}{\log m}\right).
  \]
\end{enumerate}
\end{theorem}
\begin{proof}
由 Chebyshev 函数
\(
\vartheta(x):=\sum_{p\le x}\log p
\)
与 primorial 恒等式，
\[
\log P_k=\vartheta(p_k).
\]
由素数定理的等价形式（\cite{Apostol1976}），有
\[
\vartheta(x)\sim x,\qquad p_k\sim k\log k,
\]
故
\[
\log P_k\sim k\log k.
\]
约束 \(\log_2 P_k\ge B_m(p)\) 等价于 \(\log P_k\ge S_m(p)\)，因此阈值轴数满足
\[
r_m(p)\log r_m(p)=S_m(p)\bigl(1+o(1)\bigr).
\]
设 \(u=\log r_m(p)\)，则 \(ue^u=S_m(p)(1+o(1))\)，从而
\[
u=W(S_m(p))(1+o(1)),
\qquad
r_m(p)=\frac{S_m(p)}{W(S_m(p))}(1+o(1)),
\]
得到第一式。

再由推论 \ref{cor:pom-oracle-critical-window-bit-inversion}，
\[
S_m(p)=a\,m+b\,\sqrt m+o(\sqrt m).
\]
并且 \(W(x)=\log x-\log\log x+O(1)\)（\(x\to\infty\)），故
\[
\frac{1}{W(S_m(p))}
=
\frac{1}{\log m}
+
O\!\left(\frac{\log\log m}{(\log m)^2}\right).
\]
将其与 \(S_m(p)\) 的二阶展开相乘即得
\[
r_m(p)
=
\frac{a}{\log m}\,m
+
\frac{b}{\log m}\,\sqrt m
+
O\!\left(\frac{m\log\log m}{(\log m)^2}\right)
+
o\!\left(\frac{\sqrt m}{\log m}\right).
\]
证毕。
\end{proof}

\begin{corollary}[证据预算阈值到 prime-register 轴数阶的 Lambert 反演]\label{cor:conclusion-prime-budget-axis-lambert}
设
\[
B_m(\varepsilon):=\exp\!\bigl(m\gamma_m(\varepsilon)\bigr),\qquad
r_m(\varepsilon):=\min\{r\ge 1:\ P_r>2B_m(\varepsilon)\}.
\]
若 \(m\gamma_m(\varepsilon)\to\infty\)，则
\[
r_m(\varepsilon)
=
\frac{m\gamma_m(\varepsilon)+\log 2}
{W\!\bigl(m\gamma_m(\varepsilon)+\log 2\bigr)}
\bigl(1+o(1)\bigr).
\]
进一步，在 \(\log(m\gamma_m(\varepsilon))\sim\log m\) 的常见标度区间下，
\[
r_m(\varepsilon)
=
\frac{m\gamma_m(\varepsilon)}{\log m}\bigl(1+o(1)\bigr).
\]
\end{corollary}
\begin{proof}
条件 \(P_r>2B_m(\varepsilon)\) 等价于
\[
\log P_r>m\gamma_m(\varepsilon)+\log 2.
\]
由定理 \ref{thm:conclusion-primorial-axis-second-order} 证明中的同一 primorial 渐近，
\(
\log P_r\sim r\log r
\)。
因此
\[
r_m(\varepsilon)\log r_m(\varepsilon)
=
\bigl(m\gamma_m(\varepsilon)+\log 2\bigr)\bigl(1+o(1)\bigr),
\]
Lambert \(W\) 反演得到第一式。若 \(\log(m\gamma_m(\varepsilon))\sim\log m\)，则
\(
W(m\gamma_m(\varepsilon)+\log 2)\sim\log m
\)，从而得第二式。
\end{proof}

\begin{theorem}[Gödel 整数化的对数维度税]\label{thm:conclusion-godel-logarithmic-dimension-tax}
设 prime-register
\[
\mathbf P:=\NN^{(\mathbb P)},\qquad
\mathcal N(r):=\prod_{p\in\mathbb P}p^{r(p)}.
\]
固定常数 \(A\ge 1\)。对每个 \(T\in\NN\)，考虑满足
\[
\mathrm{supp}(r)\subseteq\{p_1,\dots,p_T\},
\qquad
1\le r(p_t)\le A\ \ (1\le t\le T)
\]
的寄存器。则
\[
(1+o(1))\,T\log_2 T
\le
\log_2\mathcal N(r)
\le
(A+o(1))\,T\log_2 T
\qquad (T\to\infty).
\]
特别地，若轨迹长度 \(T=\Theta(m)\)，则其单整数 Gödel 编码位长必为 \(\Theta(m\log m)\)。
\end{theorem}
\begin{proof}
由定义
\[
\log \mathcal N(r)=\sum_{t=1}^{T}r(p_t)\log p_t.
\]
利用 \(1\le r(p_t)\le A\) 得
\[
\sum_{t=1}^{T}\log p_t
\le
\log \mathcal N(r)
\le
A\sum_{t=1}^{T}\log p_t,
\]
即
\[
\log P_T\le \log\mathcal N(r)\le A\log P_T.
\]
再由 \(\log P_T=\vartheta(p_T)\sim p_T\sim T\log T\)（\cite{Apostol1976}）得到
\[
(1+o(1))\,T\log T
\le
\log\mathcal N(r)
\le
(A+o(1))\,T\log T.
\]
两边同除以 \(\log 2\) 即得结论。
\end{proof}

\begin{proposition}[可交换投影子类中的椭圆--算法等价判据]\label{prop:conclusion-prime-register-ellipse-complete-equivalence}
\leanverified{paper\_conclusion\_prime\_register\_ellipse\_complete\_equivalence}
在 \(\mathbf P=\NN^{(\mathbb P)}\) 上定义
\[
\Psi:\mathbf P\to \mathrm{SL}(2,\RR),\qquad
\Psi(r):=\mathrm{diag}\!\bigl(\mathcal N(r),\mathcal N(r)^{-1}\bigr),
\]
并记
\[
E(r):=\Psi(r)(S^1)\subset\RR^2.
\]
则：
\begin{enumerate}
  \item \(\Psi\) 为幺半群同态（\(\Psi(r+s)=\Psi(r)\Psi(s)\)）；
  \item 映射 \(r\mapsto E(r)\) 为单射；
  \item \(\{E(r):r\in\mathbf P\}\) 恰为椭圆族
  \[
  \left\{
  \mathrm{diag}(n,n^{-1})(S^1):\ n\in\NN_{\ge 1}
  \right\},
  \]
  且该族在算子复合下与 \((\mathbf P,+)\) 同构。
\end{enumerate}
\end{proposition}
\begin{proof}
第一条由
\(
\mathcal N(r+s)=\mathcal N(r)\mathcal N(s)
\)
直接成立。第二条中，若 \(E(r)=E(s)\)，则其长半轴分别为 \(\mathcal N(r)\) 与 \(\mathcal N(s)\)，故 \(\mathcal N(r)=\mathcal N(s)\)。由算术基本定理，素数指数向量唯一，故 \(r=s\)。

对第三条，任取 \(n\in\NN_{\ge1}\)，令 \(r=v_p(n)\) 为其素数指数向量，则 \(\mathcal N(r)=n\) 且 \(E(r)=\mathrm{diag}(n,n^{-1})(S^1)\)。反向包含由定义即得。复合律
\[
\mathrm{diag}(n,n^{-1})\mathrm{diag}(m,m^{-1})
=
\mathrm{diag}(nm,(nm)^{-1})
\]
与 \(r\mapsto \mathcal N(r)\) 的乘法对应关系一致，故同构成立。
\end{proof}

\begin{proposition}[Big--Witt 谱指纹的哥德尔化可乘性]\label{prop:conclusion-bigwitt-godel-multiplicative}
设 \(T\) 为有限维（或迹类）线性算子，其动力学 \(\zeta\)-函数写为
\[
\zeta_T(z)
=\exp\!\Bigl(\sum_{n\ge 1}\frac{\Tr(T^n)}{n}z^n\Bigr)
=\prod_{n\ge 1}(1-z^n)^{-p_n(T)},
\]
其中 \(p_n(T)\in\ZZ\) 为 primitive/Witt 坐标。
记 \(\mathfrak q_n\) 为第 \(n\) 个素数，定义截断哥德尔指纹
\[
\mathcal G_N(T):=\prod_{n=1}^{N}\mathfrak q_n^{\,p_n(T)}\in\QQ_{>0}.
\]
则：
\begin{enumerate}
  \item \textbf{可恢复性}：\(\mathcal G_N(T)\) 的素因子指数唯一确定
  \((p_1(T),\dots,p_N(T))\)。
  \item \textbf{可乘性}：若 \(T=T_1\oplus T_2\)，则
  \[
  \mathcal G_N(T)=\mathcal G_N(T_1)\cdot \mathcal G_N(T_2).
  \]
  \item \textbf{谱--质数对偶}：截断 Witt 坐标的加法结构经 \(\mathcal G_N\) 严格转换为 \(\QQ_{>0}\) 的乘法结构。
\end{enumerate}
\end{proposition}
\begin{proof}
对有理数的唯一素因子分解给出
\[
v_{\mathfrak q_n}\!\bigl(\mathcal G_N(T)\bigr)=p_n(T)\qquad(1\le n\le N),
\]
故可恢复性成立。

若 \(T=T_1\oplus T_2\)，则
\[
\Tr(T^n)=\Tr(T_1^n)+\Tr(T_2^n).
\]
Witt 坐标由 ghost 坐标 \(\Tr(T^n)\) 经 Möbius/Witt 变换得到，故保持加法：
\[
p_n(T)=p_n(T_1)+p_n(T_2).
\]
代入定义即得
\[
\mathcal G_N(T)
=\prod_{n\le N}\mathfrak q_n^{p_n(T_1)+p_n(T_2)}
=\mathcal G_N(T_1)\mathcal G_N(T_2).
\]
第三条即前两条的结构表述。证毕。
\end{proof}

\begin{theorem}[原子 Big--Witt 编译的精确谱预算]\label{thm:conclusion-bigwitt-exact-spectral-budget}
\leanverified{paper\_conclusion\_bigwitt\_exact\_spectral\_budget}
设
\[
\mathcal D=\{(\beta_j,n_j,m_j)\}_{j=1}^{N}
\]
为有限原子 Euler 数据，并记其 Fredholm--Witt 目标函数为
\[
\zeta_{\mathcal D}(r):=
\prod_{j=1}^{N}(1-\beta_j r^{n_j})^{-m_j}.
\]
定义其精确谱预算
\[
\mathfrak r(\mathcal D):=
\min\Bigl\{
\operatorname{rank}_{\neq 0}(L):
\det(I-rL)^{-1}=\zeta_{\mathcal D}(r)
\Bigr\},
\]
其中 \(\operatorname{rank}_{\neq 0}(L)\) 表示 \(L\) 的非零谱代数重数总和。则成立严格恒等式
\[
\mathfrak r(\mathcal D)=\sum_{j=1}^{N}m_jn_j.
\]
并且在原子 Big--Witt 半环上有
\[
\mathfrak r(\mathcal D\oplus \mathcal E)=\mathfrak r(\mathcal D)+\mathfrak r(\mathcal E),
\qquad
\mathfrak r(\mathcal D\odot \mathcal E)=\mathfrak r(\mathcal D)\,\mathfrak r(\mathcal E).
\]
\end{theorem}
\begin{proof}
存在性已由主稿的构造性反向截面给出：对每个原子 \((\beta,n,m)\)，可取 \(m\) 个长度为 \(n\) 的循环块精确实现
\[
(1-\beta r^n)^{-m},
\]
故有限数据 \(\mathcal D\) 可由这些循环块的直和实现，其非零谱总代数重数恰为
\[
\sum_{j=1}^{N}m_jn_j.
\]
于是
\[
\mathfrak r(\mathcal D)\le \sum_{j=1}^{N}m_jn_j.
\]

反过来，若某个有限维或迹类算子 \(L\) 满足
\[
\det(I-rL)^{-1}=\zeta_{\mathcal D}(r),
\]
则 \(\det(I-rL)\) 的非平凡因子总次数不超过 \(\operatorname{rank}_{\neq 0}(L)\)。而
\[
\zeta_{\mathcal D}(r)^{-1}
=
\prod_{j=1}^{N}(1-\beta_j r^{n_j})^{m_j}
\]
的非平凡因子总次数恰为
\[
\sum_{j=1}^{N}m_jn_j.
\]
因此必有
\[
\operatorname{rank}_{\neq 0}(L)\ge \sum_{j=1}^{N}m_jn_j.
\]
两侧合并即得
\[
\mathfrak r(\mathcal D)=\sum_{j=1}^{N}m_jn_j.
\]

对直和 \(\mathcal D\oplus\mathcal E\)，目标函数相乘，对应实现为循环块直和，故非零谱总代数重数相加，从而
\[
\mathfrak r(\mathcal D\oplus\mathcal E)=\mathfrak r(\mathcal D)+\mathfrak r(\mathcal E).
\]

对 Witt 乘法 \(\mathcal D\odot\mathcal E\)，原子层上由主稿的循环块张量积恒等式知：长度 \(n\) 与长度 \(n'\) 的循环块张量后分解为
\[
\gcd(n,n')
\]
个长度
\[
\operatorname{lcm}(n,n')
\]
的循环块，因而其非零谱总代数重数变为
\[
\gcd(n,n')\operatorname{lcm}(n,n')=nn'.
\]
这正是两原子预算的乘积。再按双线性把该关系扩展到任意有限原子数据，即得
\[
\mathfrak r(\mathcal D\odot\mathcal E)=\mathfrak r(\mathcal D)\,\mathfrak r(\mathcal E).
\]
证毕。
\end{proof}

\begin{corollary}[\(p\)-典型 Frobenius--Verschiebung 的预算守恒与放大律]\label{cor:conclusion-bigwitt-typical-fv-budget-law}
\leanverified{paper\_conclusion\_bigwitt\_typical\_fv\_budget\_law}
在定理 \ref{thm:conclusion-bigwitt-exact-spectral-budget} 的记号下，对任意 \(p\)-典型原子数据 \(\mathcal D\)，成立
\[
\mathfrak r(F_p\mathcal D)=\mathfrak r(\mathcal D),
\qquad
\mathfrak r(V_p\mathcal D)=p\,\mathfrak r(\mathcal D),
\]
从而
\[
\mathfrak r(F_pV_p\mathcal D)=p\,\mathfrak r(\mathcal D).
\]
\end{corollary}
\begin{proof}
主稿的 \(p\)-典型循环块恒等式表明：Frobenius 把长度 \(p^{k+1}\) 的块送到 \(p\) 个长度 \(p^k\) 的块之直和，故非零谱总代数重数保持不变；Verschiebung 则把长度整体乘以 \(p\)，故预算整体乘以 \(p\)。于是
\[
\mathfrak r(F_p\mathcal D)=\mathfrak r(\mathcal D),
\qquad
\mathfrak r(V_p\mathcal D)=p\,\mathfrak r(\mathcal D).
\]
再应用 \(F_pV_p\) 即得最后一式。证毕。
\end{proof}

\begin{theorem}[预算可判定而函数恒等不可判定的二分律]\label{thm:conclusion-bigwitt-budget-decidable-identity-undecidable}
\leanverified{paper\_conclusion\_bigwitt\_budget\_decidable\_identity\_undecidable}
原子 Big--Witt 数据的最小精确实现预算是闭式可计算的；然而对应 Fredholm \(\zeta\) 界面的解析恒等性判定在可计算数据流层面不可判定。更准确地说：
\begin{enumerate}
  \item 对任意有限原子数据 \(\mathcal D\)，最小预算 \(\mathfrak r(\mathcal D)\) 可由
  \[
  \mathfrak r(\mathcal D)=\sum_{j=1}^{N}m_jn_j
  \]
  在有限步内计算；
  \item 存在可计算原子数据流 \(\mathcal D_0,\mathcal D_1\)，使判定
  \[
  \det(I-rL_{\mathcal D_0})^{-1}\equiv \det(I-rL_{\mathcal D_1})^{-1}
  \]
  是否在某圆盘内恒等成立，是不可判定问题。
\end{enumerate}
因此，该编译体系的大小复杂度完全刚性，而函数同一性却跨越到图灵不可判定边界。
\end{theorem}
\begin{proof}
第一点即定理 \ref{thm:conclusion-bigwitt-exact-spectral-budget} 的直接结论。

第二点正是主稿关于 Fredholm \(\zeta\) 界面等同性的不可判定性定理：存在可计算的原子数据流，使相应 Fredholm 目标函数在解析意义下是否恒等，不能由任何统一算法判定。于是最小实现预算虽可闭式计算，但解析界面的函数恒等性却仍落在不可判定边界之外。证毕。
\end{proof}

\begin{theorem}[共振审计的可判定层与 Fredholm 等同性的不可判定层]\label{thm:conclusion-resonance-decidable-fredholm-identity-undecidable-layer}
在本文所考虑的可计算原子 Euler 数据流及其有限状态同步核编译像的类中，存在如下严格的逻辑分层：
\begin{enumerate}
  \item “是否存在非平凡共振通道”是可判定的，并且可在确定性多项式时间内完成；
  \item “两份数据流所编译出的 Fredholm \(\zeta\) 界面是否在某圆盘上恒等”是不可判定的。
\end{enumerate}
因此，谱审计在本文框架中具有一条尖锐的可判定性边界：局部共振障碍可有限终止，而全局解析等同性不可有限终止。
\end{theorem}
\begin{proof}
对第 \(1\) 条，结论 \ref{con:xi-mod-obstruction-cyclotomic-equivalence} 与定理 \ref{thm:conclusion75-ihara-cyclotomic-factor-classification} 说明：存在非平凡共振通道，当且仅当某个扭曲行列式含 cyclotomic 因子。另一方面，结论 \ref{con:xi-resonant-character-snf-computable} 与命题 \ref{prop:gm-smith-compute-obstructions} 给出确定性算法，可在时间
\[
\mathrm{poly}(|V|,|E|,\log W)
\]
内输出全部障碍不变量、共振角色群的不变因子分解及生成角色表。故“是否存在非平凡共振通道”属于确定性多项式时间可判定层。

对第 \(2\) 条，定理 \ref{thm:conclusion-bigwitt-budget-decidable-identity-undecidable} 的第二点已经断言：存在可计算的原子 Euler 数据流，使相应 Fredholm \(\zeta\) 界面在某圆盘内是否恒等这一问题不可判定。两者并置，即得所述严格分层。证毕。
\end{proof}

\begin{corollary}[不存在 Fredholm--\texorpdfstring{$\zeta$}{zeta} 等同性的有限正规形]\label{cor:conclusion-no-finite-normal-form-for-fredholm-zeta-identity}
不存在任何可计算映射
\[
\mathcal N:\{\text{可计算原子 Euler 数据流}\}\longrightarrow \Sigma^{<\infty}
\]
到某个有限字母表 \(\Sigma\) 的有限字符串集合，满足
\[
\mathcal N(\mathcal D_0)=\mathcal N(\mathcal D_1)
\quad\Longleftrightarrow\quad
\zeta_{\mathcal D_0}\equiv \zeta_{\mathcal D_1}
\ \text{于某圆盘内恒等}.
\]
换言之，不存在对 Fredholm--\(\zeta\) 等同性既正确又完备的有限正规形、有限哈希或有限证书包。
\leanverified{paper\_conclusion\_no\_finite\_normal\_form\_for\_fredholm\_zeta\_identity}
\end{corollary}
\begin{proof}
反设存在这样的 \(\mathcal N\)。则对任意两份可计算原子 Euler 数据流 \(\mathcal D_0,\mathcal D_1\)，只需计算两个有限字符串并检验其是否相等，便可决定
\[
\zeta_{\mathcal D_0}\equiv \zeta_{\mathcal D_1}
\]
是否在某圆盘内恒等成立。这将把该解析等同性问题化为一个可判定问题，与定理 \ref{thm:conclusion-resonance-decidable-fredholm-identity-undecidable-layer} 的第二条矛盾。故此类有限正规形不存在。证毕。
\end{proof}


\begin{conjecture}[本质素数谱给出记录轴的精确最小支撑]\label{conj:conclusion-essential-prime-spectrum-exact-minimal-support}
设
\[
\pi_y:C\to\PP^1_y
\]
为任意绝对不可约模型在稳定反演语义下所诱导的有限覆盖，并记
\[
S_{\mathrm{ess}}(\pi_y)
\]
为定义 \ref{def:conclusion-essential-prime-spectrum} 的本质素数谱。则稳定可逆外置所需的记录轴极小支撑应当恰等于该本质素数谱；更精确地：
\begin{enumerate}
  \item 对每个
  \[
  p\notin S_{\mathrm{ess}}(\pi_y),
  \]
  其局部反演信息都可由有限标签集压缩，并与全部 \(p^k\)-级 Hensel 提升相容；
  \item 若在稳定反演制度中省略任意一个
  \[
  p\in S_{\mathrm{ess}}(\pi_y),
  \]
  则不可能保持全局单值可逆。
\end{enumerate}
换言之，记录轴的真正极小支撑不是“所有出现过的素数”，而恰是本质素数谱本身；可见相位预算与零维坏素数账本之间的和式，应当以 \(S_{\mathrm{ess}}(\pi_y)\) 为精确边界。
\end{conjecture}

\begin{conjecture}[共振梯与本质素数谱的耦合]\label{conj:conclusion-resonance-ladder-essential-prime-spectrum-coupling}
设
\[
u\longmapsto C_\varphi(u)
\]
为黄金共振梯，\(\pi_y:C\to\PP^1_y\) 为任意在稳定反演语义下讨论的绝对不可约代数覆盖，并记其本质素数谱为
\[
S_{\mathrm{ess}}(\pi_y).
\]
则应存在一个只依赖于 \(\pi_y\) 的有限素数集 \(S_\pi\) 满足：
\begin{enumerate}
  \item 对每个 \(p\notin S_\pi\)，共振梯的有限截断误差与该方向上的 Hensel 提升稳定相容；
  \item 对每个 \(p\in S_\pi\)，这种兼容性失败，且必须由外置 prime-ledger 补偿；
  \item \(S_\pi=S_{\mathrm{ess}}(\pi_y)\)。
\end{enumerate}
换言之，本质素数谱应当恰好等于共振梯发生算术分裂的最小素数支撑。
\end{conjecture}

\begin{conjecture}[本质素数谱与局部分支循环型的完备性]\label{conj:conclusion-essential-spectrum-completeness}
在固定“投影--外置--稳定反演”范式下，考虑同类代数覆盖
\(
\pi_y:C\to\PP^1
\)。
设其本质素数谱由定义 \ref{def:conclusion-essential-prime-spectrum} 给出。
若两覆盖满足
\[
S_{\mathrm{ess}}(\pi_{y,1})=S_{\mathrm{ess}}(\pi_{y,2}),
\]
且在各素数方向上的局部分支循环型（惯性群循环分解）一致，则两者在稳定可逆外置语义下同构（允许有限层 sheet 置换与局部整模型等价变换）。
\end{conjecture}

\begin{remark}[与判别式判据的接口定位]\label{rem:conclusion-essential-prime-spectrum-interface}
上述结论使用的输入仅为判别式、分支碰撞与 \(p\)-进提升稳定性三项可审计对象；与既有良约化判据相兼容，但输出目标不是单纯约化类型，而是“外置记录轴是否可被有限化”的最小性判定。
该差异使本节结果可直接接入本文的素数寄存器容量预算与算法外置语义。
\end{remark}
